\section{Dynamic control allocation for plasma shape control}
\label{sec:allocation}
In \Cref{sec:tokamaks_intro}, the second application scenario of this work has been introduced. In particular, \Cref{sec:tcv} illustrated the Tokamak à Configuration Variable device, and \Cref{sec:scd} corroborated it with a thorough description of the Système de Contrôle Distribué that controls TCV. That description clarified how the distributed control system of TCV allows multiple diagnostics, algorithms, estimators, and controller to coexist and interact within a single control architecture. Thanks to this design, the experimental capabilities of TCV have been extended since the days of the original implementation of the SCD. One of the latest innovations in this context has been the introduction of a \emph{shape controller} \cite{mele_design_2024}, \emph{i.e.}, a subsystem that acts as a reference governor and PID controller for the coil power supplies in order to track a desired plasma shape. Its implementation currently works in coordination with other SCD subsystem, in particular those that generate feedback measurements concerning the equilibrium reconstruction, that are used to evaluate the control law in real-time.

Even if the current design of the TCV shape controller might seem simple, it has been quite the challenge for the development team, and its integration and tuning have turned out to be quite delicate tasks. The end result is a system that works: given a reference plasma shape and after some tuning, it is able to control the plasma shape within the limits of the actuators and the diagnostics. Recalling what has been said in \Cref{sec:tcv}, \emph{i.e.}, the discussion of the redundant coil configuration that constitutes probably the most peculiar feature of TCV itself, one question arises: how does the shape controller take advantage of this redundancy to improve the performance of the coil current control? Up to this stage, it does not. The PID formulation of the shape controller, thoroughly described in \cite{mele_design_2024}, does not account for the actuator redundancy offered by the TCV model, which could in turn be exploited to improve the performance of the shape controller itself. Performance indicators to be considered in this context are the tracking error of the plasma shape, as well as the overall power consumption of the coils or even the saturation limits of the power supplies. In practical experiments, these have turned out to be quite important aspects to consider when working on TCV.

The aforementioned question is the starting point for the work presented in this section, which concerns a specific subject of control theory that is receiving increasing interest in recent years: that of dynamic control allocation. Techniques belonging to this subject are used to distribute the control effort among a set of redundant actuators in order to optimize a given performance index, thus following an \emph{allocation policy}. What distinguishes \emph{dynamic} control allocation from similar techniques is that, in the former, the allocation policy is a control law itself that is determined through a control design process and evaluated in real-time. Subsystems that are typically involved in a dynamic control allocation scheme are denoted as \emph{dynamic allocators} or simply as \emph{allocators}. The peculiarity of dynamic allocators is that they are designed as \emph{add-ons} to a preexisting control system, \emph{i.e.}, their design does not require a complete redesign of the control system, which could yield significant complications when optimization objectives are included in the problem. Instead, given the models of the plant and the control system, their design can be carried out \emph{a posteriori}, and then they can be integrated into the control system as additional components properly connected to the control loop. The art of dynamic control allocation, therefore, lies in the design of these additional components with the goal in mind of not only optimizing the control action in terms of some criteria of practical relevance, but also of ensuring that the original control objectives are still met, possibly without any alteration of the desired output behavior.

Regarding TCV, a preliminary study has been carried out in \cite{tenaglia_dynamic_2024}, in which an early implementation of a dynamic coil current allocator has been integrated in the MATLAB code suite of TCV. Its performance has been validated by means of simulations on real shot data, which showed the feasibility of this approach. This section presents the work that has been carried out in this direction. First, a review of the current state of the art about dynamic control allocation is presented, together with a set of formulations that are of interest for the TCV case. Each of these has been considered, leading to some independent theoretical results that are reported in the following. Then, the design of a novel dynamic control allocator for TCV is presented, which is based on the latest advancements on the subject. At the end of the section, experimental results are discussed, which show the performance of the proposed novel allocator in a real scenario and constitute the first working implementation of its kind in the literature.
